\documentclass[11pt]{amsart}

\usepackage[margin=1.15in]{geometry}
\usepackage{amsmath,amssymb,amsthm,mathtools}
\usepackage{enumitem}
\usepackage{microtype}
\usepackage[colorlinks=true,linkcolor=blue,citecolor=blue,urlcolor=blue]{hyperref}
\usepackage[nameinlink,capitalize]{cleveref}

\theoremstyle{plain}
\newtheorem{theorem}{Theorem}[section]
\newtheorem{proposition}[theorem]{Proposition}
\newtheorem{lemma}[theorem]{Lemma}
\newtheorem{corollary}[theorem]{Corollary}

\theoremstyle{definition}
\newtheorem{definition}[theorem]{Definition}
\newtheorem{assumption}[theorem]{Assumption}
\newtheorem{hypothesis}[theorem]{Hypothesis}

\theoremstyle{remark}
\newtheorem{remark}[theorem]{Remark}

\numberwithin{equation}{section}

\DeclareMathOperator{\diver}{div}

\title[Structured ancient branch exclusions]
{Structured Ancient Navier--Stokes Solutions and Lyapunov Branch Exclusions}
\author{Guillem Duran Ballester}
\email{guillem@fragile.tech}
\date{}

\subjclass[2020]{35Q30, 35B40, 35B44, 35B45, 76D05}
\keywords{Navier--Stokes equations, ancient solutions, Type I blow-up, axisymmetric flows, Lyapunov functionals, Burgers vortices}

\begin{document}

\begin{abstract}
We study structured bounded ancient solutions of the three-dimensional
Navier--Stokes equations arising in the Type I blow-up reduction.  We separate
the structured branch into closed Liouville subclasses, perturbative Lyapunov
subclasses, model-vortex subclasses, and a remaining bounded-swirl obstruction.

For the unforced equations, several axisymmetric branches are closed by known
Liouville theorems: no swirl, pointwise scale-invariant control, finite
\(L_t^\infty L_x^p\) control of the swirl variable
\(\Gamma=ru_\theta\), and axial periodicity with bounded \(\Gamma\).  When such
branches occur as normalized Type I ancient limits, the Galilean-trivial
conclusions of the external Liouville theorems reduce to the zero solution and
contradict the blow-up normalization.

We also record two Lyapunov-based exclusions.  First, Gallay--Wayne's
small-data decay theorem for the three-dimensional rescaled vorticity equation
yields a coercive Lyapunov functional excluding perturbative weighted-vorticity
ancient trajectories.  Second, Gallay--Maekawa's stability theorem for the
strained Burgers-vortex model yields a profile-relative Lyapunov functional and
ancient rigidity inside a small weighted tube around the Burgers-vortex family.

Finally, for the open nonperiodic bounded-swirl axisymmetric branch, we prove
the exact scalar \(\Gamma\)-defect identity and the anisotropic shell-error
decomposition.  These identities give a coercive shell-control criterion under
which \(\Gamma\) must vanish, reducing the solution to a no-swirl branch.  The
paper concludes with a branch ledger specifying which structured alternatives are
closed, which are conditional, and which analytic inputs remain.
\end{abstract}

\maketitle

\section{Introduction}

\subsection{Structured ancient branches in the Type I program}

The Type I blow-up reduction in Paper I \cite{DuranPaperI} associates an
admissible Type I singularity with a normalized nonzero Seregin ancient limit.
In physical variables this limit is represented by a smooth ancient solution
\((U,P)\) of
\begin{equation}\label{eq:NS}
   \partial_tU+(U\cdot\nabla)U+\nabla P=\Delta U,
   \qquad \nabla\cdot U=0,
   \qquad t<0,
\end{equation}
satisfying the Type I ancient bound
\begin{equation}\label{eq:typeI}
   \sup_{t<0}\sqrt{-t}\,
   \|U(t)\|_{L^\infty(\mathbb R^3)}<\infty .
\end{equation}
The extraction also gives a nonvanishing normalization.  Thus every theorem in
this paper is used in the Type I program through the same pattern:
\[
   \text{structured ancient branch}
   \quad\Longrightarrow\quad
   U\equiv0
   \quad\Longrightarrow\quad
   \text{contradiction with normalization}.
\]

The present paper treats only those structured alternatives for which a
specific closure mechanism is available or for which a precise conditional
criterion can be stated.  The branches considered are:
\[
   \begin{gathered}
   \text{axisymmetric Liouville branches},\qquad
   \text{perturbative weighted-vorticity branch},\\
   \text{Burgers-vortex model branch},\qquad
   \text{bounded-swirl shell branch}.
   \end{gathered}
\]
It does not claim that every structured or decaying ancient solution is closed.
The output is a branch ledger for the structured/decay part of the residual
criterion of Paper I.

\subsection{Closed axisymmetric Liouville branches}

Several unforced axisymmetric branches are already closed by external
Liouville theorems.  Koch--Nadirashvili--Seregin--\v{S}ver\'ak close the no-swirl
branch and the pointwise scale-invariant branch \(|u|\le C/r\).  The
swirl-bearing finite-\(L^p\) branch
\[
   \Gamma=ru_\theta\in L_t^\infty L_x^p,\qquad 1\le p<\infty,
\]
is closed by Lei--Zhang--Zhao.  The \(z\)-periodic bounded-\(\Gamma\) branch is
closed by Lei--Ren--Zhang.

The only extra step needed for Type I ancient limits is to convert
Galilean-triviality into zero in the blow-up frame.  We time-shift away from
the possible singular time \(t=0\), apply the bounded ancient Liouville
theorem, patch the resulting constants on backward intervals, and use
\eqref{eq:typeI}: if \(U(x,t)\equiv c\), then
\[
   \sup_{t<0}\sqrt{-t}\,\|U(t)\|_\infty
   =
   \sup_{t<0}\sqrt{-t}\,|c|
\]
is finite only when \(c=0\).

\subsection{Lyapunov branches}

Two branch exclusions below are Lyapunov arguments.  The first is the
Gallay--Wayne perturbative weighted-vorticity branch.  If \(w\) solves the
rescaled three-dimensional vorticity equation and remains in the
Gallay--Wayne small ball of \(L^2_\sigma(m)\), \(m>7/2\), then the functional
\begin{equation}\label{eq:intro-weighted-lyap}
   \mathcal L_m(w_0)
   =
   \sup_{\tau\ge0}e^{2\tau}\|\Phi_\tau w_0\|_m^2
\end{equation}
is coercive and contracts along the forward semiflow.  An ancient trajectory
inside the small ball is therefore zero after sending the starting time to
\(-\infty\).

The second Lyapunov branch is the strained Burgers-vortex model.  The theorem
proved here is a theorem for
\[
   \partial_\tau\Omega+(U\cdot\nabla)\Omega-(\Omega\cdot\nabla)U
   =
   \Delta\Omega-(Ax\cdot\nabla)\Omega+A\Omega,
   \qquad
   A=\operatorname{diag}\left(-\frac12,-\frac12,1\right),
\]
not directly for the unforced Navier--Stokes equations.  Gallay--Maekawa's
stability theorem gives a profile-relative Lyapunov functional on a fixed
circulation fiber and forces ancient trajectories in a small weighted tube to
be exactly stationary Burgers vortices.  Transferring that conclusion to an
unforced branch requires the model-reduction assumption stated later.

\subsection{The bounded-swirl obstruction}

The remaining axisymmetric obstruction is the nonperiodic bounded-swirl branch
\[
   \Gamma=ru_\theta\in L^\infty
      \bigl(\mathbb R^3\times(-\infty,0)\bigr),
\]
with no finite-\(L^p\) control, no axial periodicity, and no decay hypothesis.
For a smooth axisymmetric solution, the scalar \(\Gamma\) satisfies
\begin{equation}\label{eq:intro-gamma}
   \partial_t\Gamma+b\cdot\nabla\Gamma+\frac2r\partial_r\Gamma
   =
   \Delta\Gamma,
   \qquad
   b=u_re_r+u_ze_z .
\end{equation}
The corresponding weighted identity is
\begin{equation}\label{eq:intro-gamma-identity}
   \frac12\frac{d}{dt}\int\Gamma^2\phi\,dx
   +
   \int|\nabla\Gamma|^2\phi\,dx
   =
   \frac12\int\Gamma^2
   \left(
      \partial_t\phi+\Delta\phi+b\cdot\nabla\phi+\frac2r\partial_r\phi
   \right)\,dx .
\end{equation}
For anisotropic shell weights
\[
   \phi_{R,Z,\zeta}(r,z,t)=\lambda_R(r)\eta_Z(z-\zeta(t)),
\]
the error density splits exactly into radial shell transport, axial drift
mismatch, and cutoff curvature.  The radial diffusion operator is
\(\partial_r^2+3r^{-1}\partial_r\), the radial Laplacian in four dimensions.
The paper proves that a coercive shell-control hypothesis forces
\(\Gamma\equiv0\).  It does not prove that every bounded nonperiodic
\(\Gamma\) satisfies this shell-control hypothesis.

\subsection{Main results}

The first result closes the unforced axisymmetric branches relevant to the
Type I program.

\begin{theorem}[Closed axisymmetric ancient branches]
\label{thm:closed-axisymmetric-branches}
Let \(U\) be a smooth Type I ancient solution of \eqref{eq:NS} satisfying
\eqref{eq:typeI}.  Assume \(U\) is a normalized ancient limit produced by a
Type I blow-up sequence.  If \(U\) belongs to any of the following
axisymmetric branches, then \(U\equiv0\):
\begin{enumerate}[label=\textup{(\roman*)}]
\item axisymmetric without swirl;
\item axisymmetric and satisfying a pointwise scale-invariant bound
\[
   |U(x,t)|\le \frac{C}{r};
\]
\item axisymmetric with
\[
   \Gamma=rU_\theta\in L_t^\infty L_x^p
   \bigl(\mathbb R^3\times(-\infty,0)\bigr),
   \qquad 1\le p<\infty;
\]
\item axisymmetric, periodic in the physical axial variable \(z\), and with
\[
   \Gamma=rU_\theta\in L^\infty
   \bigl(\mathbb R^3\times(-\infty,0)\bigr).
\]
\end{enumerate}
Consequently none of these branches can occur as a nonzero normalized Type I
ancient limit.
\end{theorem}

\begin{theorem}[Perturbative weighted-vorticity Lyapunov exclusion]
\label{thm:weighted-vorticity-exclusion-intro}
Let \(m>\frac72\).  Let \(w(\tau)\) be an ancient mild solution of the
Gallay--Wayne rescaled vorticity equation
\begin{equation}\label{eq:weighted-vorticity-intro}
   \partial_\tau w
   =
   \Lambda w-(v\cdot\nabla)w+(w\cdot\nabla)v,
   \qquad \nabla\cdot w=0,
\end{equation}
with \(v\) recovered by the three-dimensional Biot--Savart law.  Assume
\[
   \sup_{\tau\in\mathbb R}\|w(\tau)\|_{L^2_\sigma(m)}
   \le \rho_m,
\]
where \(\rho_m\) is the Gallay--Wayne small-data radius.  Then
\[
   w\equiv0.
\]
\end{theorem}

Here ``ancient mild'' means that every forward time-shift of \(w\) is the
Gallay--Wayne mild solution launched from its value at the starting time.

\begin{theorem}[Ancient rigidity in the Burgers-vortex tube]
\label{thm:burgers-vortex-tube-intro}
Let \(m>2\).  Let
\(\Omega:\mathbb R\to\mathcal X_{\mathrm{full}}(m)\) be an admissible ancient
mild solution of the strained Burgers-vortex vorticity equation
\begin{equation}\label{eq:burgers-intro}
   \partial_\tau\Omega+(U\cdot\nabla)\Omega-(\Omega\cdot\nabla)U
   =
   \Delta\Omega-(Ax\cdot\nabla)\Omega+A\Omega,
   \qquad
   A=\operatorname{diag}\left(-\frac12,-\frac12,1\right).
\end{equation}
Assume \(\Omega\) remains in the Gallay--Maekawa small weighted tube around the
Burgers-vortex family, where \(\bar\alpha\) denotes the circulation coordinate
of \(\Omega\):
\[
   \sup_{\tau\in\mathbb R}
   \|\Omega(\tau)-\bar\alpha G\|_{\mathcal X_{\mathrm{mod}}(m)}
   \le
   \delta_{\bar\alpha,m}.
\]
Then
\[
   \Omega(\tau)\equiv\bar\alpha G
   \qquad\text{for all }\tau\in\mathbb R.
\]
\end{theorem}

This theorem is a theorem for the strained Burgers-vortex model.  It becomes a
statement about unforced Navier--Stokes ancient limits only after an additional
model-reduction hypothesis, stated as \Cref{ass:bv-reduction}.  The term
``ancient mild'' includes compatibility with the forward Gallay--Maekawa
semiflow after every time shift, as made explicit in \Cref{def:bv-tube}.

\begin{theorem}[Conditional scalar flattening]
\label{thm:scalar-flattening-intro}
Let \(u\) be a smooth bounded ancient mild axisymmetric solution of the
unforced Navier--Stokes equations with bounded swirl variable
\[
   \Gamma=ru_\theta\in L^\infty.
\]
Assume the coercive shell-control hypotheses stated in
\Cref{ass:shell-control}.  Then
\[
   \Gamma\equiv0
   \quad\text{on }\mathbb R^3\times(-\infty,0).
\]
Consequently, if the resulting no-swirl solution lies in a closed no-swirl
Liouville branch, then \(u\) is Galilean-trivial; if in addition \(u\) is a
Type I ancient limit, then \(u\equiv0\).
\end{theorem}

\begin{theorem}[Structured branch input for the residual criterion]
\label{thm:structured-input-paperI}
Let \(V\) be a normalized Seregin ancient limit generated by an admissible
Type I source, and let \(U\) be its physical ancient representative.  Suppose
one of the following precise alternatives holds: \(U\) belongs to one of the
closed structured branches listed in \Cref{thm:closed-axisymmetric-branches};
the associated rescaled vorticity belongs to the perturbative
weighted-vorticity branch closed by
\Cref{thm:weighted-vorticity-exclusion-intro}; the associated model variable
belongs to the Burgers-vortex tube under the model-reduction hypothesis
\Cref{ass:bv-reduction}; or \(U\) belongs to the bounded-swirl branch under
\Cref{ass:shell-control}.  In the first, second, and fourth alternatives the
corresponding normalized branch is zero and hence cannot occur as a nonzero
residual-limit alternative in Paper I.  In the Burgers-vortex alternative, the
conclusion is the model rigidity statement
\(\mathcal T[U](\tau)\equiv\bar\alpha G\); the model theorem alone makes no
unforced Navier--Stokes zero-solution assertion.
\end{theorem}

The theorem applies only to the closed subclasses and to the conditional
subclasses for which the stated hypotheses hold.

\subsection{Relation to Papers I and II}

Paper I \cite{DuranPaperI} reduces admissible Type I singularities to
normalized nonzero Seregin ancient limits and formulates a residual-class
criterion.  Paper II \cite{DuranPaperII} treats the uniformly \(L^3\)-tight
branch.  The present paper treats structured branches.  Its output is a
structured-branch ledger: several structured alternatives are excluded
unconditionally by external Liouville theorems, some perturbative branches are
excluded by Lyapunov arguments, and the remaining bounded-swirl axisymmetric
branch is reduced to a precise shell-control criterion.  Paper III does not
depend on Paper II.

\subsection{Branch ledger}

\begin{center}
\begin{tabular}{p{0.36\textwidth}p{0.18\textwidth}p{0.34\textwidth}}
\hline
Branch & Status & Input \\
\hline
axisymmetric no swirl & closed & KNSS \\
pointwise \(|u|\le C/r\) & closed & KNSS \\
\(\Gamma\in L_t^\infty L_x^p\), \(1\le p<\infty\) & closed & LZZ \\
\(z\)-periodic bounded \(\Gamma\) & closed & LRZ \\
controlled swirl \(\rho^\gamma V_\theta\) & conditional & fixed-exponent hypothesis \\
broad weighted velocity decay & conditional & fixed-exponent hypothesis \\
small weighted vorticity & closed & Gallay--Wayne \\
Burgers-vortex tube & model closed & Gallay--Maekawa \\
unforced Burgers-vortex transfer & conditional & model-reduction assumption \\
bounded nonperiodic \(\Gamma\in L^\infty\) & open; conditional criterion & shell-control criterion \\
bounded \(\Gamma\) plus shell control & closed & scalar flattening plus no-swirl \\
\hline
\end{tabular}
\end{center}

\subsection{Organization}

\Cref{sec:typeI} records the Type I ancient setup, self-similar variables, the
structure-preservation facts, the time-shift lemma, and the branch classes.
\Cref{sec:axis-inputs} fixes axisymmetric notation and the external Liouville
inputs.  \Cref{sec:closed-axisymmetric} proves the closed unforced
axisymmetric branch theorem.  \Cref{sec:conditional} records the controlled
swirl and broad weighted velocity-decay alternatives as conditional hypotheses.
\Cref{sec:weighted-vorticity} proves the Gallay--Wayne Lyapunov exclusion.
\Cref{sec:burgers} proves the Burgers-vortex model rigidity.  \Cref{sec:gamma}
derives the scalar \(\Gamma\)-identity, shell-error decomposition, and
conditional scalar flattening theorem.  \Cref{sec:ledger} gives the final
branch ledger and the interface with Paper I.

\section{Ancient Type I limits, self-similar variables, and structured classes}
\label{sec:typeI}

\subsection{Type I ancient solutions}

\begin{definition}[Type I ancient solution]
\label{def:typeI-ancient}
A smooth mild solution \(U\) of \eqref{eq:NS} on
\(\mathbb R^3\times(-\infty,0)\) is called Type I ancient if
\eqref{eq:typeI} holds.  Mild is understood in the standard velocity
formulation \cite{Kato1984,Sohr2001}: for every \(-\infty<s<t<0\),
\begin{equation}\label{eq:mild-formula}
   U(t)=e^{(t-s)\Delta}U(s)
   -\int_s^t e^{(t-\sigma)\Delta}\mathbb P\nabla\cdot
      (U(\sigma)\otimes U(\sigma))\,d\sigma
\end{equation}
in distributions, where \(\mathbb P\) is the Leray projector.
\end{definition}

\subsection{Normalized ancient limits}

\begin{definition}[Normalized Type I ancient limit]
\label{def:normalized-typeI-limit}
A Type I ancient solution \(U\) is called a normalized Type I ancient limit if
it arises from the blow-up extraction of Paper I and satisfies the
corresponding nonvanishing normalization.  Equivalently, its centered
representative \(V\) is a normalized Seregin ancient limit in the sense of
Paper I.
\end{definition}

\begin{lemma}[Nonzero normalization]
\label{lem:normalized-nonzero}
A normalized Type I ancient limit is not identically zero.
\end{lemma}

\begin{proof}
This is the nonvanishing conclusion of the Type I extraction theorem in
Paper I.  The present paper uses only the consequence \(U\not\equiv0\), not
the particular normalization functional.
\end{proof}

\subsection{Backward self-similar variables}

For \(t<0\), define
\begin{equation}\label{eq:self-similar}
   V(y,\tau)=\sqrt{-t}\,U(x,t),\qquad
   Q(y,\tau)=(-t)P(x,t),\qquad
   y=\frac{x}{\sqrt{-t}},\qquad
   \tau=-\log(-t).
\end{equation}
Equivalently,
\[
   x=e^{-\tau/2}y,\qquad t=-e^{-\tau}.
\]

\begin{lemma}[Self-similar equation]\label{lem:self-similar-equation}
Let \((U,P)\) be a smooth solution of \eqref{eq:NS}, and define \((V,Q)\) by
\eqref{eq:self-similar}.  Then \(V\) is defined on
\(\mathbb R^3\times\mathbb R\) and satisfies
\begin{equation}\label{eq:centered}
   \partial_\tau V+(V\cdot\nabla)V+\nabla Q
   =
   \Delta V-\frac12(V+y\cdot\nabla V),
   \qquad \nabla\cdot V=0.
\end{equation}
Moreover \eqref{eq:typeI} is equivalent to
\[
   \sup_{\tau\in\mathbb R}
   \|V(\tau)\|_{L^\infty(\mathbb R^3)}<\infty .
\]
\end{lemma}

\begin{proof}
Let \(a=\sqrt{-t}=e^{-\tau/2}\).  Then \(V=aU\), \(y=x/a\), and
\[
   \partial_t\tau=\frac1{-t}=\frac1{a^2},
   \qquad
   \partial_t y=\frac{y}{2a^2},
   \qquad
   \partial_t(a^{-1})=\frac1{2a^3}.
\]
Since \(U=a^{-1}V\), differentiating at fixed \(x\) gives
\[
   \partial_tU
   =
   a^{-3}
   \left(
      \partial_\tau V+\frac12 y\cdot\nabla V+\frac12V
   \right).
\]
Spatial derivatives scale as
\[
   \nabla_xU=a^{-2}\nabla_yV,\qquad
   \Delta_xU=a^{-3}\Delta_yV,
\]
and the pressure scaling gives
\[
   \nabla_xP=a^{-3}\nabla_yQ.
\]
Thus
\[
   (U\cdot\nabla_x)U=a^{-3}(V\cdot\nabla_y)V.
\]
Multiplying \eqref{eq:NS} by \(a^3\) yields \eqref{eq:centered}; the
divergence equation follows from
\(\nabla_x\cdot U=a^{-2}\nabla_y\cdot V\).  The spatial map
\(y\mapsto x=e^{-\tau/2}y\) is a bijection for each \(\tau\), and
\[
   \|V(\tau)\|_\infty
   =
   \sqrt{-t}\,\|U(t)\|_\infty .
\]
Taking suprema over \(t<0\), equivalently \(\tau\in\mathbb R\), gives the
last assertion.
\end{proof}

\subsection{Preservation of axisymmetry and swirl variables}

For axisymmetric flows we write
\[
   x=(r\cos\theta,r\sin\theta,z),\qquad
   U=U_re_r+U_\theta e_\theta+U_ze_z,
\]
and in centered variables
\[
   y=(\rho\cos\theta,\rho\sin\theta,Z),\qquad
   V=V_\rho e_\rho+V_\theta e_\theta+V_Ze_Z.
\]
The scalar swirl variables are
\[
   \Gamma_U(r,z,t)=rU_\theta(r,z,t),\qquad
   \Gamma_V(\rho,Z,\tau)=\rho V_\theta(\rho,Z,\tau).
\]

\begin{lemma}[Preservation of axisymmetric structure]
\label{lem:structure-preserved}
Axisymmetry and absence of swirl are preserved under
\eqref{eq:self-similar}.  If \(U\) is axisymmetric, then
\[
   \Gamma_V(\rho,Z,\tau)=\Gamma_U(r,z,t),
   \qquad
   (r,z)=\sqrt{-t}\,(\rho,Z).
\]
Fixed physical \(z\)-periodicity is preserved by physical time shifts, but
under the centered change of variables it becomes a \(\tau\)-dependent period
in the \(Z\)-coordinate.
\end{lemma}

\begin{proof}
The self-similar map is a scalar spatial dilation and a scalar multiplication
of the velocity, so it commutes with rotations around the symmetry axis.
The cylindrical basis is unchanged by radial dilation, and
\[
   V_\theta(\rho,Z,\tau)
   =
   \sqrt{-t}\,U_\theta(r,z,t),
   \qquad
   r=\sqrt{-t}\rho,\quad z=\sqrt{-t}Z .
\]
Thus \(U_\theta\equiv0\) if and only if \(V_\theta\equiv0\), and
\[
   \Gamma_V=\rho V_\theta
   =
   \rho\sqrt{-t}\,U_\theta
   =
   rU_\theta
   =
   \Gamma_U .
\]
If \(U\) has physical axial period \(L\), then \(V\) has period
\(L/\sqrt{-t}=Le^{\tau/2}\) in \(Z\), not a fixed period.
\end{proof}

\subsection{Time shifts away from the singular time}

\begin{lemma}[Time shifts away from the singular time]
\label{lem:time-shift}
Let \(U\) be a smooth Type I ancient solution and fix \(T<0\).  Then
\[
   W_T(x,s)=U(x,T+s),\qquad \Pi_T(x,s)=P(x,T+s),\qquad s<0,
\]
is a smooth bounded mild ancient solution of \eqref{eq:NS} on
\(\mathbb R^3\times(-\infty,0)\).  Moreover the axisymmetric, no-swirl,
\(L^p\)-swirl, bounded-\(\Gamma\), and physical \(z\)-periodic properties are
preserved.
\end{lemma}

\begin{proof}
Time translation preserves the differential equation because
\(\partial_sW_T(x,s)=\partial_tU(x,T+s)\) and all spatial derivatives are
unchanged.  For the mild formulation, apply \eqref{eq:mild-formula} to \(U\)
between the physical times \(T+s<T+t<0\):
\[
   U(T+t)=e^{(t-s)\Delta}U(T+s)
   -\int_{T+s}^{T+t}e^{(T+t-\eta)\Delta}\mathbb P\nabla\cdot
      (U(\eta)\otimes U(\eta))\,d\eta .
\]
With \(\eta=T+\sigma\), this becomes the mild formula for \(W_T\) on the
interval \((s,t)\).  For \(s<0\), \(T+s<T<0\), so the Type I bound gives
\[
   \|W_T(s)\|_\infty
   =
   \|U(T+s)\|_\infty
   \le
   \frac1{\sqrt{-T}}
   \sup_{\eta<0}\sqrt{-\eta}\,\|U(\eta)\|_\infty .
\]
The structural assertions follow because only the time variable has been
translated.  In particular
\[
   \Gamma_{W_T}(r,z,s)=\Gamma_U(r,z,T+s),
\]
so \(L_s^\infty L_x^p\) and \(L^\infty\) bounds for \(\Gamma\) restrict to the
shifted solution, and a fixed physical axial period remains the same.
\end{proof}

\subsection{Structured branch classes}

\begin{definition}[Closed axisymmetric branch classes]
\label{def:closed-axisymmetric-classes}
A Type I ancient solution \(U\) belongs to:
\begin{enumerate}[label=\textup{(\roman*)}]
\item \(\mathfrak C_{\mathrm{noswirl}}\) if it is axisymmetric and
\(U_\theta\equiv0\);
\item \(\mathfrak C_{1/r}\) if it is axisymmetric and
\[
   |U(x,t)|\le C/r \qquad(r>0);
\]
\item \(\mathfrak C_{\Gamma,p}\), \(1\le p<\infty\), if it is
axisymmetric and
\[
   \Gamma=rU_\theta\in L_t^\infty L_x^p
   \bigl(\mathbb R^3\times(-\infty,0)\bigr);
\]
\item \(\mathfrak C_{\mathrm{per},\Gamma}\) if it is axisymmetric, periodic
in the physical axial variable \(z\), and \(\Gamma\in L^\infty\).
\end{enumerate}
\end{definition}

\begin{definition}[Conditional structured branch classes]
\label{def:conditional-structured-classes}
We also record the following conditional branches:
\begin{enumerate}[label=\textup{(\roman*)}]
\item the controlled-swirl branch \(\mathfrak C_{\mathrm{cs}}(\gamma)\);
\item the broad weighted velocity-decay branch \(\mathfrak C_{\mathrm{wd}}(m)\);
\item the perturbative weighted-vorticity branch \(\mathfrak C_{\mathrm{GW}}(m)\);
\item the Burgers-vortex reduction branch \(\mathfrak C_{\mathrm{BV}}(m)\);
\item the bounded-swirl shell-control branch \(\mathfrak C_{\mathrm{shell}}\).
\end{enumerate}
The first two require fixed-exponent Liouville hypotheses.  The Burgers-vortex
branch requires a model-reduction assumption when it is imported into
unforced Navier--Stokes.  The shell-control branch requires
\Cref{ass:shell-control}.
\end{definition}

\section{Axisymmetric notation and external Liouville inputs}
\label{sec:axis-inputs}

\subsection{Cylindrical variables and the swirl scalar}

For a generic axisymmetric solution \(u\), write
\[
   u=u_re_r+u_\theta e_\theta+u_ze_z,\qquad
   \Gamma=ru_\theta,\qquad
   b=u_re_r+u_ze_z .
\]
When \(\Gamma\) is placed in \(L^p(\mathbb R^3)\), it is regarded as the
corresponding function of \(x=(r\cos\theta,r\sin\theta,z)\), independent of
\(\theta\):
\[
   \|\Gamma(t)\|_{L^p(\mathbb R^3)}^p
   =
   2\pi\int_{\mathbb R}\int_0^\infty
      |\Gamma(r,z,t)|^p r\,dr\,dz,\qquad 1\le p<\infty .
\]

\subsection{Constant and Galilean-trivial solutions}

\begin{definition}[Galilean-trivial]
A velocity field is Galilean-trivial if, after a Galilean transform, it is a
constant vector field.  In the Type I blow-up frame below we use only the
following consequence: when an external Liouville theorem gives a constant
velocity field in that frame, the Type I bound forces the constant to be zero.
\end{definition}

\begin{lemma}[Spatially constant mild solutions]\label{lem:spatial-constant-mild}
Let \(W\) be a smooth mild solution of \eqref{eq:NS} on
\(\mathbb R^3\times(-\infty,0)\).  If \(W(x,t)=b(t)\) for all \(x\) and \(t\),
then \(b\) is constant in time.
\end{lemma}

\begin{proof}
Insert \(W(x,t)=b(t)\) into \eqref{eq:mild-formula}.  The heat semigroup
leaves constant vector fields fixed, and
\(\nabla\cdot(b(\sigma)\otimes b(\sigma))=0\).  Thus
\(b(t)=b(s)\) for every \(s<t<0\).
\end{proof}

\begin{lemma}[Patching constants on backward intervals]
\label{lem:patch-constants}
Let \(U:\mathbb R^3\times(-\infty,0)\to\mathbb R^3\) be a function.  Assume
that for every \(T<0\) there exists \(c_T\in\mathbb R^3\) such that
\[
   U(x,t)=c_T\qquad(x\in\mathbb R^3,\ t<T).
\]
Then there is a single \(c\in\mathbb R^3\) such that
\[
   U(x,t)=c\qquad(x\in\mathbb R^3,\ t<0).
\]
\end{lemma}

\begin{proof}
If \(T_1<T_2<0\), the two constants agree on the nonempty interval
\((-\infty,T_1)\).  Hence all \(c_T\) are equal.  For any \(t<0\), choose
\(T\) with \(t<T<0\).
\end{proof}

\begin{lemma}[Constant Type I ancient solutions vanish]
\label{lem:galilean-typeI-zero}
Let \(U\) be a Type I ancient solution.  If \(U(x,t)\equiv c\), then
\(c=0\).
\end{lemma}

\begin{proof}
The Type I bound gives
\[
   \sqrt{-t}\,|c|
   \le
   \sup_{s<0}\sqrt{-s}\,\|U(s)\|_\infty<\infty
   \qquad(t<0).
\]
Letting \(t\to-\infty\) gives \(c=0\).
\end{proof}

\subsection{External Liouville inputs}

The following theorems are used as external PDE inputs.

\begin{theorem}[KNSS no-swirl Liouville theorem]
\label{thm:knss-noswirl}
Every bounded ancient mild axisymmetric no-swirl solution of the unforced
Navier--Stokes equations is Galilean-trivial.  In the form used below, the
solution is a spatially constant axial drift.
\end{theorem}

\begin{proof}
This is the no-swirl Liouville theorem of
Koch--Nadirashvili--Seregin--\v{S}ver\'ak \cite{KNSS2009}, stated in the form needed
for the branch exclusion.
\end{proof}

\begin{theorem}[KNSS pointwise scale-invariant branch]
\label{thm:knss-pointwise}
Every bounded ancient mild axisymmetric solution satisfying
\[
   |u(x,t)|\le \frac{C}{r}
   \qquad(r>0,\ t<0)
\]
is Galilean-trivial.  In the form used below it is identically zero.
\end{theorem}

\begin{proof}
This is the pointwise branch theorem of \cite{KNSS2009}, stated in the form
used here.
\end{proof}

\begin{theorem}[Lei--Zhang--Zhao]
\label{thm:lzz}
Let \(u\) be a bounded ancient mild axisymmetric solution.  If
\[
   \Gamma=ru_\theta\in L_t^\infty L_x^p
   \bigl(\mathbb R^3\times(-\infty,0)\bigr),
   \qquad 1\le p<\infty,
\]
then \(u\) is a constant vector field.
\end{theorem}

\begin{proof}
This is the ancient axisymmetric \(L^p\)-swirl Liouville theorem of
Lei--Zhang--Zhao \cite{LZZ}.
\end{proof}

\begin{theorem}[Lei--Ren--Zhang]
\label{thm:lrz}
Let \(u\) be a bounded ancient mild axisymmetric solution, periodic in \(z\),
with
\[
   \Gamma=ru_\theta\in L^\infty
   \bigl(\mathbb R^3\times(-\infty,0)\bigr).
\]
Then \(u\) is a constant axial vector field.
\end{theorem}

\begin{proof}
This is the periodic bounded-swirl theorem of Lei--Ren--Zhang
\cite{LRZ}.  Their statement is on \(\mathbb R^2\times\mathbb T^1\); a
physical \(z\)-periodic whole-space solution descends to the quotient, and a
Navier--Stokes scaling normalizes the period.
\end{proof}

\section{Closed unforced axisymmetric branches}
\label{sec:closed-axisymmetric}

\subsection{No swirl}

\begin{proposition}[No-swirl Type I exclusion]\label{prop:noswirl-typeI}
Let \(U\) be a smooth Type I ancient solution in
\(\mathfrak C_{\mathrm{noswirl}}\).  Then \(U\equiv0\).
\end{proposition}

\begin{proof}
Fix \(T<0\) and set \(W_T(x,s)=U(x,T+s)\).  By \Cref{lem:time-shift},
\(W_T\) is a bounded ancient mild axisymmetric no-swirl solution.  Applying
\Cref{thm:knss-noswirl} gives that \(W_T\) is a spatially constant axial
drift, say \(W_T(x,s)=b_T(s)e_z\).  By
\Cref{lem:spatial-constant-mild}, \(b_T(s)\) is constant in \(s\).  Thus
\(U(x,t)=c_T\) for \(t<T\).  Since \(T<0\) is arbitrary,
\Cref{lem:patch-constants} gives \(U(x,t)=c\) for all \(t<0\).  Finally,
\Cref{lem:galilean-typeI-zero} gives \(c=0\).
\end{proof}

\subsection{Pointwise scale-invariant control}

\begin{proposition}[Pointwise \(1/r\) Type I exclusion]\label{prop:pointwise-typeI}
Let \(U\) be a smooth Type I ancient solution in \(\mathfrak C_{1/r}\).  Then
\(U\equiv0\).
\end{proposition}

\begin{proof}
Fix \(T<0\).  The shifted solution \(W_T(x,s)=U(x,T+s)\) is bounded, ancient,
axisymmetric, and satisfies the same pointwise bound \(|W_T|\le C/r\).
\Cref{thm:knss-pointwise} gives \(W_T\equiv0\).  Hence \(U\equiv0\) on
\((-\infty,T)\).  Since \(T<0\) is arbitrary, \(U\equiv0\).
\end{proof}

\subsection{\texorpdfstring{Finite \(L^p\) control of \(\Gamma\)}
{Finite Lp control of Gamma}}

\begin{proposition}[Finite-\(L^p\) swirl Type I exclusion]
\label{prop:lp-swirl-typeI}
Let \(U\) be a smooth Type I ancient solution in
\(\mathfrak C_{\Gamma,p}\) for some \(1\le p<\infty\).  Then \(U\equiv0\).
\end{proposition}

\begin{proof}
Fix \(T<0\).  By \Cref{lem:time-shift}, \(W_T(x,s)=U(x,T+s)\) is a bounded
ancient mild axisymmetric solution and
\[
   \Gamma_{W_T}\in L_s^\infty L_x^p
   \bigl(\mathbb R^3\times(-\infty,0)\bigr).
\]
\Cref{thm:lzz} gives that \(W_T\) is a constant vector field.  Patch these
constants in \(T\) by \Cref{lem:patch-constants} and then use
\Cref{lem:galilean-typeI-zero}.
\end{proof}

\subsection{Periodic bounded swirl}

\begin{proposition}[Periodic bounded-swirl Type I exclusion]
\label{prop:periodic-swirl-typeI}
Let \(U\) be a smooth Type I ancient solution in
\(\mathfrak C_{\mathrm{per},\Gamma}\).  Then \(U\equiv0\).
\end{proposition}

\begin{proof}
Fix \(T<0\).  By \Cref{lem:time-shift}, the shifted solution \(W_T\) is
bounded, ancient, axisymmetric, physically \(z\)-periodic, and has bounded
\(\Gamma\).  \Cref{thm:lrz} gives that \(W_T\) is a constant axial vector
field.  The constants patch as \(T\) varies, and the Type I bound eliminates
the resulting constant by \Cref{lem:galilean-typeI-zero}.
\end{proof}

\subsection{Proof of the closed branch theorem}

\begin{proof}[Proof of \Cref{thm:closed-axisymmetric-branches}]
The four alternatives are exactly the hypotheses of
\Cref{prop:noswirl-typeI,prop:pointwise-typeI,prop:lp-swirl-typeI,prop:periodic-swirl-typeI}.
Each proposition gives \(U\equiv0\).  This contradicts
\Cref{lem:normalized-nonzero} when \(U\) is a normalized Type I ancient limit.
\end{proof}

\section{Conditional controlled-swirl and weighted-decay hypotheses}
\label{sec:conditional}

\subsection{Controlled swirl in self-similar variables}

\begin{definition}[Controlled-swirl branch]
\label{def:controlled-swirl}
Fix \(\gamma>0\).  A centered ancient solution \(V\) belongs to the
controlled-swirl branch \(\mathfrak C_{\mathrm{cs}}(\gamma)\) if \(V\) is
axisymmetric and
\[
   \sup_{\tau\in\mathbb R}\sup_{\rho>0,\ Z\in\mathbb R}
   \rho^\gamma |V_\theta(\rho,Z,\tau)|<\infty .
\]
\end{definition}

\begin{hypothesis}[Fixed-exponent controlled-swirl Liouville theorem]
\label{hyp:controlled-swirl-liouville}
For the chosen exponent \(\gamma>0\), every normalized Type I ancient limit
whose centered representative lies in
\(\mathfrak C_{\mathrm{cs}}(\gamma)\) is zero.
\end{hypothesis}

\begin{proposition}[Conditional controlled-swirl exclusion]
\label{prop:controlled-swirl-conditional}
If \Cref{hyp:controlled-swirl-liouville} holds for a fixed \(\gamma>0\), then
no nonzero normalized Type I ancient limit belongs to
\(\mathfrak C_{\mathrm{cs}}(\gamma)\).
\end{proposition}

\begin{proof}
This is exactly the hypothesis combined with the nonzero normalization
\Cref{lem:normalized-nonzero}.
\end{proof}

\subsection{Broad weighted velocity-decay branch}

\begin{definition}[Broad weighted velocity-decay branch]
\label{def:weighted-velocity-branch}
Fix \(m>0\).  A centered ancient solution \(V\) belongs to
\(\mathfrak C_{\mathrm{wd}}(m)\) if
\[
   \sup_{\tau\in\mathbb R}
   \int_{\mathbb R^3}(1+|y|^2)^m |V(y,\tau)|^2\,dy<\infty .
\]
\end{definition}

\begin{hypothesis}[Fixed-exponent weighted velocity Liouville theorem]
\label{hyp:weighted-velocity-liouville}
For the chosen exponent \(m>0\), every normalized Type I ancient limit whose
centered representative lies in \(\mathfrak C_{\mathrm{wd}}(m)\) is zero.
\end{hypothesis}

\begin{proposition}[Conditional weighted velocity exclusion]
\label{prop:weighted-velocity-conditional}
If \Cref{hyp:weighted-velocity-liouville} holds for a fixed \(m>0\), then no
nonzero normalized Type I ancient limit belongs to
\(\mathfrak C_{\mathrm{wd}}(m)\).
\end{proposition}

\begin{proof}
This is \Cref{hyp:weighted-velocity-liouville} plus
\Cref{lem:normalized-nonzero}.
\end{proof}

\begin{remark}[Weighted velocity versus weighted vorticity]
The broad weighted velocity branch is not closed by the Gallay--Wayne theorem
below.  Gallay--Wayne applies only to ancient trajectories whose rescaled
vorticity remains in a sufficiently small ball of \(L^2_\sigma(m)\).  Thus
\Cref{hyp:weighted-velocity-liouville} is a separate fixed-exponent Liouville
input unless an additional reduction from broad weighted velocity decay to
perturbative weighted vorticity is supplied.
\end{remark}

\section{Perturbative weighted-vorticity Lyapunov exclusion}
\label{sec:weighted-vorticity}

\subsection{The rescaled vorticity equation}

In this section \(\tau\) is the forward self-similar time used in the
Gallay--Wayne vorticity formulation.  The ancient trajectories are indexed by
all \(\tau\in\mathbb R\), and the forward semiflow acts by increasing
\(\tau\).

Let \(w=\nabla\times v\).  The rescaled vorticity equation is
\begin{equation}\label{eq:weighted-vorticity}
   \partial_\tau w
   =
   \Lambda w-(v\cdot\nabla)w+(w\cdot\nabla)v,
   \qquad \nabla\cdot w=0,
\end{equation}
where
\begin{equation}\label{eq:lambda}
   \Lambda w=\Delta w+\frac12\xi\cdot\nabla w+w,
\end{equation}
and \(v\) is recovered from \(w\) by the three-dimensional Biot--Savart law.
For \(m\ge0\), set
\[
   L^2_\sigma(m)
   =
   \left\{
      f\in L^2(\mathbb R^3;\mathbb R^3):
      \nabla\cdot f=0,\quad \|f\|_m<\infty
   \right\},
\]
where
\[
   \|f\|_m^2
   =
   \int_{\mathbb R^3}(1+|\xi|)^{2m}|f(\xi)|^2\,d\xi .
\]

\subsection{Gallay--Wayne small-data decay}

\begin{theorem}[Gallay--Wayne small-data decay]\label{thm:gallay}
Let \(0<\mu\le1\) and \(m>2\mu+\frac12\).  There exist constants
\(\rho_m>0\) and \(C_m\ge1\) such that if \(\|w_0\|_m\le\rho_m\), then
\eqref{eq:weighted-vorticity} has a unique global mild solution
\(w(\tau)=\Phi_\tau w_0\) in \(C([0,\infty);L^2_\sigma(m))\) and
\[
   \|\Phi_\tau w_0\|_m
   \le
   C_m e^{-\mu\tau}\|w_0\|_m,\qquad \tau\ge0.
\]
In particular, if \(m>\frac72\), one may use \(\mu=1\).
\end{theorem}

\begin{proof}
This is the Gallay--Wayne weighted small-data theorem for the rescaled
three-dimensional vorticity equation \cite{GallayWayne2002}.  The only
conclusions used below are existence, uniqueness, the semiflow property, and
the displayed exponential decay.
\end{proof}

\subsection{The Lyapunov functional}

\begin{definition}[Perturbative weighted-vorticity branch]
\label{def:weighted-vorticity-branch}
Fix \(m>\frac72\), and let \(\rho_m\) be the radius in \Cref{thm:gallay}.
The branch \(\mathfrak C_{\mathrm{GW}}(m)\) consists of all
\[
   w\in C(\mathbb R;L^2_\sigma(m))
\]
which solve \eqref{eq:weighted-vorticity} for all \(\tau\in\mathbb R\) and
satisfy
\[
   \sup_{\tau\in\mathbb R}\|w(\tau)\|_m\le\rho_m .
\]
Solving for all \(\tau\) means that every time shift is compatible with the
Gallay--Wayne forward mild semiflow.
\end{definition}

\begin{proposition}[Semiflow Lyapunov functional]\label{prop:weighted-lyap}
Fix \(m>\frac72\).  For \(\|w_0\|_m\le\rho_m\), define
\begin{equation}\label{eq:weighted-lyap}
   \mathcal L_m(w_0)
   =
   \sup_{\tau\ge0}e^{2\tau}\|\Phi_\tau w_0\|_m^2 .
\end{equation}
Then \(\mathcal L_m\) is well-defined and satisfies
\begin{equation}\label{eq:weighted-coercive}
   \|w_0\|_m^2
   \le
   \mathcal L_m(w_0)
   \le
   C_m^2\|w_0\|_m^2 .
\end{equation}
Moreover, for every \(s\ge0\),
\begin{equation}\label{eq:weighted-contract}
   \mathcal L_m(\Phi_s w_0)
   \le
   e^{-2s}\mathcal L_m(w_0).
\end{equation}
\end{proposition}

\begin{proof}
The Gallay--Wayne estimate with \(\mu=1\) gives
\[
   e^{2\tau}\|\Phi_\tau w_0\|_m^2
   \le C_m^2\|w_0\|_m^2,
\]
which proves finiteness and the upper bound.  The lower bound is the value at
\(\tau=0\).  By uniqueness,
\[
   \Phi_\tau(\Phi_s w_0)=\Phi_{\tau+s}w_0 .
\]
Hence
\[
   \mathcal L_m(\Phi_s w_0)
   =
   \sup_{\tau\ge0}e^{2\tau}\|\Phi_{\tau+s}w_0\|_m^2
   =
   e^{-2s}\sup_{\sigma\ge s}e^{2\sigma}\|\Phi_\sigma w_0\|_m^2
   \le e^{-2s}\mathcal L_m(w_0).
\]
\end{proof}

\subsection{Ancient rigidity in the small weighted ball}

\begin{proof}[Proof of \Cref{thm:weighted-vorticity-exclusion-intro}]
Let \(w\) be an ancient solution in the small ball and fix
\(\tau_0\in\mathbb R\).  For \(s<\tau_0\), the shifted function
\[
   W_s(\theta)=w(s+\theta),\qquad \theta\ge0,
\]
is the Gallay--Wayne mild solution with initial datum \(w(s)\).  The small-ball
hypothesis gives \(\|w(s)\|_m\le\rho_m\), so the forward solution is in the
domain of \(\Phi_\theta\), and uniqueness gives
\[
   W_s(\theta)=\Phi_\theta w(s),\qquad \theta\ge0.
\]
Choosing \(\theta=\tau_0-s\) gives
\[
   w(\tau_0)=\Phi_{\tau_0-s}w(s).
\]
Since \(\|w(\tau_0)\|_m\le\rho_m\), the Lyapunov functional is also defined at
\(w(\tau_0)\).  Applying \Cref{prop:weighted-lyap} with forward time
\(\tau_0-s\),
\[
   \mathcal L_m(w(\tau_0))
   \le
   e^{-2(\tau_0-s)}\mathcal L_m(w(s))
   \le
   C_m^2\rho_m^2e^{-2(\tau_0-s)} .
\]
Letting \(s\to-\infty\) gives \(\mathcal L_m(w(\tau_0))=0\).  The lower
coercive bound \eqref{eq:weighted-coercive} gives \(w(\tau_0)=0\).  Since
\(\tau_0\) was arbitrary, \(w\equiv0\).
\end{proof}

\begin{remark}[Relation to weighted velocity decay]
This theorem supplies a structured/decay branch exclusion only for normalized
limits whose vorticity can be placed inside the Gallay--Wayne perturbative
weighted branch.  It does not by itself exclude all weighted velocity-decay
alternatives in the branch taxonomy of Paper I.
\end{remark}

\section{The Burgers-vortex model branch}
\label{sec:burgers}

This section proves a rigidity theorem for the strained Burgers-vortex model.
It is not, by itself, a theorem for the unforced Navier--Stokes equations.  The
final corollary states the additional model-reduction assumption needed to
transfer the conclusion to an unforced ancient-solution class.

\subsection{The strained vorticity equation}

Let \(x=(x_h,x_3)\in\mathbb R^2\times\mathbb R\) and
\[
   A=\operatorname{diag}\left(-\frac12,-\frac12,1\right).
\]
The Burgers-vortex vorticity equation is
\begin{equation}\label{eq:burgers}
   \partial_\tau\Omega+(U\cdot\nabla)\Omega-(\Omega\cdot\nabla)U
   =
   \Delta\Omega-(Ax\cdot\nabla)\Omega+A\Omega,
   \qquad
   \nabla\cdot\Omega=0,
\end{equation}
with \(U\) recovered from \(\Omega\) by the Biot--Savart law.  It has the
stationary family
\begin{equation}\label{eq:burgers-family}
   \Omega=\alpha G,\qquad
   G(x)=\begin{pmatrix}0\\0\\g(x_h)\end{pmatrix},
   \qquad
   g(x_h)=\frac1{4\pi}e^{-|x_h|^2/4},
   \qquad \alpha\in\mathbb R .
\end{equation}

\subsection{Weighted horizontal spaces}

For \(m>1\), let \(L^2_h(m)\) denote the horizontal weighted space with norm
\[
   \|f\|_{L^2_h(m)}^2
   =
   \int_{\mathbb R^2}|f(x_h)|^2
   \left(1+\frac{|x_h|^2}{4m}\right)^m\,dx_h .
\]
Set
\[
   X(m)=BC(\mathbb R_{x_3};L^2_h(m)),\qquad
   X_0(m)=
   \left\{
      f\in X(m):
      \int_{\mathbb R^2}f(x_h,x_3)\,dx_h=0
      \text{ for every }x_3
   \right\},
\]
and
\[
   \mathcal X_{\mathrm{full}}(m)=X(m)^3,\qquad
   \mathcal X_{\mathrm{mod}}(m)=X(m)\times X(m)\times X_0(m).
\]

\subsection{Admissible Burgers-vortex solutions}

\begin{definition}[Admissible Burgers-vortex solutions]\label{def:admissible-bv}
A mild solution
\(\Omega:\mathbb R\to\mathcal X_{\mathrm{full}}(m)\) is called admissible if
the following operations are justified on every compact time interval:
\begin{enumerate}[label=\textup{(\roman*)}]
\item horizontal integrations by parts against constants and against \(x_h\)
are valid, with no boundary contribution at \(|x_h|=\infty\);
\item the identity
\[
   \int_{\mathbb R^2}
   \bigl[(U\cdot\nabla)\Omega-(\Omega\cdot\nabla)U\bigr]_3\,dx_h=0
\]
holds in distributions in \(x_3\) and \(\tau\);
\item the map
\[
   \tau\mapsto
   \int_{\mathbb R^2}\Omega_3(x_h,x_3,\tau)\,dx_h
\]
is absolutely continuous after testing in \(x_3\).
\end{enumerate}
Smooth rapidly decaying solutions are admissible.
\end{definition}

\subsection{Circulation coordinate and conservation}

\begin{lemma}[Circulation coordinate]\label{lem:circulation}
Let \(m>1\) and let
\(\Omega\in\mathcal X_{\mathrm{full}}(m)\) be divergence-free.  Assume that
the horizontal integration by parts
\[
   \int_{\mathbb R^2}\nabla_h\cdot\Omega_h(x_h,x_3)\,dx_h=0
\]
is valid.  Then
\[
   \alpha[\Omega](x_3)
   =
   \int_{\mathbb R^2}\Omega_3(x_h,x_3)\,dx_h
\]
is well-defined and independent of \(x_3\).  Writing the common value as
\(\alpha[\Omega]\), one has
\[
   Q\Omega:=\Omega-\alpha[\Omega]G\in\mathcal X_{\mathrm{mod}}(m).
\]
\end{lemma}

\begin{proof}
Since \(m>1\), \(L^2_h(m)\hookrightarrow L^1(\mathbb R^2)\) by
Cauchy--Schwarz:
\[
   \int_{\mathbb R^2}|f(x_h)|\,dx_h
   \le
   \|f\|_{L^2_h(m)}
   \left(
      \int_{\mathbb R^2}
      \left(1+\frac{|x_h|^2}{4m}\right)^{-m}\,dx_h
   \right)^{1/2},
\]
and the last integral is finite in two dimensions exactly because \(m>1\).
Thus \(\alpha[\Omega](x_3)\) is well-defined.  Using
\(\nabla\cdot\Omega=0\),
\[
   \partial_{x_3}\alpha[\Omega](x_3)
   =
   \int_{\mathbb R^2}\partial_{x_3}\Omega_3\,dx_h
   =
   -\int_{\mathbb R^2}\nabla_h\cdot\Omega_h\,dx_h
   =
   0 .
\]
Also \(\int_{\mathbb R^2}g(x_h)\,dx_h=1\), so subtracting
\(\alpha[\Omega]G\) removes the horizontal mean of the third component.
\end{proof}

\begin{lemma}[Conservation of circulation]\label{lem:circulation-conserved}
Let \(m>1\), and let \(\Omega(\tau)\) be an admissible mild solution of
\eqref{eq:burgers} with values in \(\mathcal X_{\mathrm{full}}(m)\).  Then
\(\alpha[\Omega(\tau)]\) is independent of \(\tau\).  If
\(\bar\alpha\) denotes this constant, then
\[
   \omega(\tau)=\Omega(\tau)-\bar\alpha G
\]
belongs to \(\mathcal X_{\mathrm{mod}}(m)\) for every \(\tau\) and solves the
fixed-circulation perturbation equation around \(\bar\alpha G\).
\end{lemma}

\begin{proof}
Integrate the third component of \eqref{eq:burgers} over \(x_h\).  The
nonlinear term is a horizontal divergence:
\[
   (U\cdot\nabla)\Omega-(\Omega\cdot\nabla)U
   =
   -\nabla\times(U\times\Omega),
\]
and the third component of a curl is a horizontal divergence.  It therefore
has zero horizontal integral by admissibility.  For the linear part,
\[
   \bigl[\Delta\Omega-(Ax\cdot\nabla)\Omega+A\Omega\bigr]_3
   =
   \Delta_h\Omega_3+\partial_{x_3}^2\Omega_3
   +\frac12x_h\cdot\nabla_h\Omega_3
   -x_3\partial_{x_3}\Omega_3+\Omega_3 .
\]
The \(\Delta_h\) term integrates to zero, and
\[
   \int_{\mathbb R^2}\frac12x_h\cdot\nabla_h\Omega_3\,dx_h
   =
   -\int_{\mathbb R^2}\Omega_3\,dx_h,
\]
which cancels the \(+\Omega_3\) term.  Hence
\[
   \frac{d}{d\tau}\alpha[\Omega(\tau)]
   =
   \partial_{x_3}^2\alpha[\Omega(\tau)]
   -x_3\partial_{x_3}\alpha[\Omega(\tau)].
\]
By \Cref{lem:circulation}, the right-hand side is zero.
\end{proof}

\subsection{Gallay--Maekawa stability}

\begin{theorem}[Gallay--Maekawa stability]\label{thm:gm}
Let \(m>2\) and \(\alpha\in\mathbb R\).  There exist constants
\(\delta_{\alpha,m}>0\) and \(C_{\alpha,m}\ge1\) such that every
divergence-free perturbation
\[
   \omega_0\in\mathcal X_{\mathrm{mod}}(m),
   \qquad
   \|\omega_0\|_{\mathcal X_{\mathrm{mod}}(m)}
   \le\delta_{\alpha,m},
\]
generates a unique global mild solution of the fixed-circulation perturbation
equation around \(\alpha G\), denoted \(\Phi_\tau^\alpha\omega_0\), and
\[
   \|\Phi_\tau^\alpha\omega_0\|_{\mathcal X_{\mathrm{mod}}(m)}
   \le
   C_{\alpha,m}e^{-\tau/2}
   \|\omega_0\|_{\mathcal X_{\mathrm{mod}}(m)}
   \qquad(\tau\ge0).
\]
\end{theorem}

\begin{proof}
This is the Gallay--Maekawa nonlinear stability theorem for Burgers vortices
in horizontal weighted spaces \cite{GallayMaekawa}.  We use only existence,
uniqueness, the fixed-circulation semiflow property, and the displayed
exponential decay.
\end{proof}

\subsection{Profile-relative Lyapunov functional}

\begin{proposition}[Profile-relative Lyapunov functional]\label{prop:bv-lyap}
Fix \(m>2\) and \(\alpha\in\mathbb R\).  On the ball
\[
   \|\omega_0\|_{\mathcal X_{\mathrm{mod}}(m)}
   \le\delta_{\alpha,m}
\]
define
\begin{equation}\label{eq:bv-lyap}
   \mathcal B_{\alpha,m}(\omega_0)
   =
   \sup_{\tau\ge0}e^\tau
   \|\Phi_\tau^\alpha\omega_0\|_{\mathcal X_{\mathrm{mod}}(m)}^2 .
\end{equation}
Then
\[
   \|\omega_0\|_{\mathcal X_{\mathrm{mod}}(m)}^2
   \le
   \mathcal B_{\alpha,m}(\omega_0)
   \le
   C_{\alpha,m}^2
   \|\omega_0\|_{\mathcal X_{\mathrm{mod}}(m)}^2,
\]
and
\[
   \mathcal B_{\alpha,m}(\Phi_s^\alpha\omega_0)
   \le
   e^{-s}\mathcal B_{\alpha,m}(\omega_0)
   \qquad(s\ge0).
\]
\end{proposition}

\begin{proof}
Gallay--Maekawa stability gives, for every \(\tau\ge0\),
\[
   e^\tau
   \|\Phi_\tau^\alpha\omega_0\|_{\mathcal X_{\mathrm{mod}}(m)}^2
   \le
   C_{\alpha,m}^2
   \|\omega_0\|_{\mathcal X_{\mathrm{mod}}(m)}^2 .
\]
Taking the supremum over \(\tau\ge0\) proves that
\(\mathcal B_{\alpha,m}\) is finite and gives the upper coercive bound.  The
lower bound follows from the value at \(\tau=0\), because
\(\Phi_0^\alpha\omega_0=\omega_0\).

For \(s\ge0\), uniqueness of the fixed-circulation perturbation equation gives
the semiflow identity
\[
   \Phi_\tau^\alpha(\Phi_s^\alpha\omega_0)
   =
   \Phi_{\tau+s}^\alpha\omega_0 .
\]
Therefore
\[
   \mathcal B_{\alpha,m}(\Phi_s^\alpha\omega_0)
   =
   \sup_{\tau\ge0}
   e^\tau
   \|\Phi_{\tau+s}^\alpha\omega_0\|_{\mathcal X_{\mathrm{mod}}(m)}^2 .
\]
With \(\sigma=\tau+s\), this is
\[
   e^{-s}\sup_{\sigma\ge s}
   e^\sigma
   \|\Phi_\sigma^\alpha\omega_0\|_{\mathcal X_{\mathrm{mod}}(m)}^2
   \le
   e^{-s}\mathcal B_{\alpha,m}(\omega_0).
\]
\end{proof}

\subsection{Ancient rigidity in the Burgers-vortex tube}

\begin{definition}[Burgers-vortex ancient tube]\label{def:bv-tube}
Fix \(m>2\).  The Burgers-vortex ancient tube consists of all admissible
ancient mild solutions \(\Omega:\mathbb R\to\mathcal X_{\mathrm{full}}(m)\)
of \eqref{eq:burgers} such that, with
\(\bar\alpha=\alpha[\Omega(\tau)]\),
\[
   \sup_{\tau\in\mathbb R}
   \|\Omega(\tau)-\bar\alpha G\|_{\mathcal X_{\mathrm{mod}}(m)}
   \le
   \delta_{\bar\alpha,m}.
\]
The term ``ancient mild'' includes compatibility with the
Gallay--Maekawa forward mild semiflow after every time shift.
\end{definition}

\begin{proof}[Proof of \Cref{thm:burgers-vortex-tube-intro}]
By \Cref{lem:circulation-conserved}, \(\bar\alpha\) is independent of time and
\[
   \omega(\tau)=\Omega(\tau)-\bar\alpha G
\]
is an ancient solution of the fixed-circulation perturbation equation in the
small ball.  Fix \(\tau_0\in\mathbb R\) and \(s<\tau_0\).  By the ancient mild
compatibility, the shifted perturbation
\[
   \omega_s(\theta)=\omega(s+\theta),\qquad \theta\ge0,
\]
is the Gallay--Maekawa mild solution launched from \(\omega(s)\).  The tube
hypothesis gives
\[
   \|\omega(s)\|_{\mathcal X_{\mathrm{mod}}(m)}
   \le\delta_{\bar\alpha,m},
\]
so uniqueness of the fixed-circulation semiflow gives
\[
   \omega(\tau_0)=\Phi_{\tau_0-s}^{\bar\alpha}\omega(s).
\]
Using \Cref{prop:bv-lyap},
\[
   \mathcal B_{\bar\alpha,m}(\omega(\tau_0))
   \le
   e^{-(\tau_0-s)}
   \mathcal B_{\bar\alpha,m}(\omega(s))
   \le
   e^{-(\tau_0-s)}
   C_{\bar\alpha,m}^2\delta_{\bar\alpha,m}^2 .
\]
Letting \(s\to-\infty\) gives
\(\mathcal B_{\bar\alpha,m}(\omega(\tau_0))=0\).  The coercive lower bound
gives \(\omega(\tau_0)=0\).  Since \(\tau_0\) was arbitrary,
\(\Omega(\tau)\equiv\bar\alpha G\).
\end{proof}

\subsection{Conditional transfer to unforced Navier--Stokes}

\begin{assumption}[Reduction to the Burgers-vortex model]\label{ass:bv-reduction}
Let \(\mathcal C\) be a class of ancient solutions of an unforced or
renormalized Navier--Stokes problem.  Assume there are \(m>2\), a map
\[
   \mathcal T:\mathcal C\to
   C_{\mathrm{loc}}(\mathbb R;\mathcal X_{\mathrm{full}}(m)),
\]
and a nonnegative defect \(d[u(\tau)]\) such that for every \(u\in\mathcal C\):
\begin{enumerate}[label=\textup{(\roman*)}]
\item \(\Omega=\mathcal T[u]\) is an admissible ancient mild solution of
\eqref{eq:burgers} in the sense of \Cref{def:admissible-bv};
\item with \(\bar\alpha=\alpha[\Omega]\),
\[
   \|\Omega(\tau)-\bar\alpha G\|_{\mathcal X_{\mathrm{mod}}(m)}
   \le C_* d[u(\tau)]
   \qquad(\tau\in\mathbb R);
\]
\item \(\sup_{\tau\in\mathbb R}d[u(\tau)]\le
\delta_{\bar\alpha,m}/C_*\).
\end{enumerate}
\end{assumption}

\begin{corollary}[Conditional Burgers-vortex model transfer]
\label{cor:bv-import}
Under \Cref{ass:bv-reduction}, every \(u\in\mathcal C\) is mapped by
\(\mathcal T\) to an exact Burgers vortex:
\[
   \mathcal T[u](\tau)\equiv\bar\alpha G .
\]
\end{corollary}

\begin{proof}
The assumptions put \(\Omega=\mathcal T[u]\) in the Burgers-vortex ancient
tube of \Cref{def:bv-tube}.  Applying
\Cref{thm:burgers-vortex-tube-intro} gives the conclusion.
\end{proof}

\section{The scalar \texorpdfstring{\(\Gamma\)}{Gamma} defect and shell-control criterion}
\label{sec:gamma}

\subsection{The equation for \texorpdfstring{\(\Gamma\)}{Gamma}}

\begin{lemma}[Swirl equation]\label{lem:gamma-equation}
Let \(u\) be a smooth axisymmetric solution of \eqref{eq:NS}.  Then
\(\Gamma=ru_\theta\) satisfies
\begin{equation}\label{eq:gamma}
   \partial_t\Gamma
   +b\cdot\nabla\Gamma
   +\frac2r\partial_r\Gamma
   =
   \Delta\Gamma,
   \qquad b=u_re_r+u_ze_z .
\end{equation}
\end{lemma}

\begin{proof}
The angular component of \eqref{eq:NS} is
\[
   \partial_tu_\theta+u_r\partial_r u_\theta+u_z\partial_z u_\theta
   +\frac{u_r}{r}u_\theta
   =
   \left(\partial_r^2+\frac1r\partial_r+\partial_z^2-\frac1{r^2}\right)
   u_\theta .
\]
Multiplying by \(r\), the transport terms become
\[
   r u_r\partial_ru_\theta+u_ru_\theta
   =
   u_r\partial_r(ru_\theta)
   =
   u_r\partial_r\Gamma,
   \qquad
   r u_z\partial_zu_\theta
   =
   u_z\partial_z\Gamma .
\]
For the radial diffusion, use \(u_\theta=\Gamma/r\).  Then
\[
   \partial_ru_\theta
   =
   \frac1r\partial_r\Gamma-\frac1{r^2}\Gamma,
   \qquad
   \partial_r^2u_\theta
   =
   \frac1r\partial_r^2\Gamma
   -\frac2{r^2}\partial_r\Gamma
   +\frac2{r^3}\Gamma .
\]
Consequently
\[
   r\left(\partial_r^2+\frac1r\partial_r-\frac1{r^2}\right)u_\theta
   =
   \partial_r^2\Gamma-\frac1r\partial_r\Gamma .
\]
The axial diffusion satisfies \(r\partial_z^2u_\theta=\partial_z^2\Gamma\).
Since
\[
   \Delta\Gamma
   =
   \partial_r^2\Gamma+\frac1r\partial_r\Gamma+\partial_z^2\Gamma,
\]
the total diffusion term after multiplication by \(r\) is
\[
   \partial_r^2\Gamma-\frac1r\partial_r\Gamma+\partial_z^2\Gamma
   =
   \Delta\Gamma-\frac2r\partial_r\Gamma .
\]
Substitution gives
\[
   \partial_t\Gamma+b\cdot\nabla\Gamma
   =
   \Delta\Gamma-\frac2r\partial_r\Gamma,
\]
which is \eqref{eq:gamma}.
\end{proof}

\subsection{The weighted scalar defect identity}

\begin{proposition}[Weighted scalar defect identity]\label{prop:gamma-identity}
Let \(u\) be a smooth axisymmetric solution of \eqref{eq:NS}.  Set
\(W=\Gamma\), and let \(\phi=\phi(r,z,t)\ge0\) be a smooth compactly supported
axisymmetric weight.  Then
\begin{equation}\label{eq:gamma-identity}
   \frac12\frac{d}{dt}\int_{\mathbb R^3}W^2\phi\,dx
   +
   \int_{\mathbb R^3}|\nabla W|^2\phi\,dx
   =
   \frac12\int_{\mathbb R^3}W^2
   \left(
      \partial_t\phi+\Delta\phi+b\cdot\nabla\phi+\frac2r\partial_r\phi
   \right)\,dx .
\end{equation}
If \(W=\Gamma-c\), the same identity holds provided \(c=0\) or \(\phi\)
vanishes near the axis.  Without this restriction, the right-hand side must be
augmented by the axis boundary term
\[
   2\pi c^2\int_{\mathbb R}\phi(0,z,t)\,dz .
\]
\end{proposition}

\begin{proof}
From \Cref{lem:gamma-equation},
\[
   \partial_tW
   =
   \Delta W-b\cdot\nabla W-\frac2r\partial_rW .
\]
Multiply by \(W\phi\) and integrate.  The time derivative gives
\[
   \int W\partial_tW\,\phi\,dx
   =
   \frac12\frac{d}{dt}\int W^2\phi\,dx
   -
   \frac12\int W^2\partial_t\phi\,dx .
\]
The Laplacian gives
\[
   \int W\Delta W\,\phi\,dx
   =
   -\int|\nabla W|^2\phi\,dx
   +
   \frac12\int W^2\Delta\phi\,dx .
\]
Because \(u\) is divergence-free and \(b\) is its meridional part,
\(\nabla\cdot b=0\) in the three-dimensional sense:
\[
   \nabla\cdot b
   =
   \frac1r\partial_r(ru_r)+\partial_zu_z
   =
   \nabla\cdot u
   =
   0.
\]
Thus, using compact support of \(\phi\),
\[
   -\int W(b\cdot\nabla W)\phi\,dx
   =
   \frac12\int W^2 b\cdot\nabla\phi\,dx .
\]
Finally, in cylindrical variables \(dx=r\,dr\,d\theta\,dz\),
\[
   -\int W\frac2r\partial_rW\,\phi\,dx
   =
   -\int_0^{2\pi}\int_{\mathbb R}\int_0^\infty
      \partial_r(W^2)\phi\,dr\,dz\,d\theta .
\]
For \(W=\Gamma\), smoothness gives \(\Gamma=O(r^2)\) at the axis, so the
boundary term at \(r=0\) vanishes.  More explicitly, after integrating on
\(\delta<r<\infty\) and letting \(\delta\downarrow0\), the boundary
contribution is
\[
   \int_0^{2\pi}\int_{\mathbb R}W(\delta,z,t)^2\phi(\delta,z,t)\,dz\,d\theta
   \longrightarrow0 .
\]
Therefore the last display equals
\[
   \int W^2\frac1r\partial_r\phi\,dx
   =
   \frac12\int W^2\frac2r\partial_r\phi\,dx .
\]
Combining the terms proves \eqref{eq:gamma-identity}.  If \(W=\Gamma-c\),
then \(W(0,z,t)=-c\).  The same integration on \(\delta<r<\infty\) leaves the
limit
\[
   \int_0^{2\pi}\int_{\mathbb R}c^2\phi(0,z,t)\,dz\,d\theta
   =
   2\pi c^2\int_{\mathbb R}\phi(0,z,t)\,dz ,
\]
which is absent exactly when \(c=0\) or \(\phi\) avoids the axis.
\end{proof}

\subsection{Anisotropic shell weights}

Let
\[
   \phi_{R,Z,\zeta}(r,z,t)=\lambda_R(r)\eta_Z(z-\zeta(t)),
\]
where \(\lambda_R,\eta_Z\) are smooth cutoffs and \(\zeta\) is a smooth axial
modulation.

\begin{proposition}[Shell-error decomposition]\label{prop:shell}
Let \(W=\Gamma\).  With the notation above,
\begin{equation}\label{eq:shell-identity}
   \frac12\frac{d}{dt}\int W^2\lambda_R\eta_Z\,dx
   +
   \int|\nabla W|^2\lambda_R\eta_Z\,dx
   =
   \frac12\int W^2\mathcal E_{R,Z,\zeta}[u]\,dx,
\end{equation}
where
\begin{equation}\label{eq:shell-error}
   \mathcal E_{R,Z,\zeta}[u]
   =
   \eta_Z\left(
      \lambda_R''+\frac3r\lambda_R'+u_r\lambda_R'
   \right)
   +
   \lambda_R\left(
      \eta_Z''+(u_z-\dot\zeta)\eta_Z'
   \right).
\end{equation}
\end{proposition}

\begin{proof}
For an axisymmetric scalar,
\[
   \Delta(\lambda_R\eta_Z)
   =
   \left(\lambda_R''+\frac1r\lambda_R'\right)\eta_Z
   +
   \lambda_R\eta_Z'' .
\]
Also
\[
   \frac2r\partial_r(\lambda_R\eta_Z)
   =
   \frac2r\lambda_R'\eta_Z,\qquad
   b\cdot\nabla(\lambda_R\eta_Z)
   =
   u_r\lambda_R'\eta_Z+u_z\lambda_R\eta_Z',
\]
and
\[
   \partial_t(\lambda_R\eta_Z)
   =
   -\dot\zeta\,\lambda_R\eta_Z'.
\]
Adding these four contributions gives \eqref{eq:shell-error}; substitution
into \Cref{prop:gamma-identity} gives \eqref{eq:shell-identity}.
\end{proof}

\begin{remark}[The radial geometry]
The differential expression
\[
   \lambda_R''+\frac3r\lambda_R'
\]
is the radial Laplacian in four dimensions.  Thus the natural radial barrier
for the \(\Gamma\)-defect is adapted to
\(\partial_r^2+3r^{-1}\partial_r\).  The shell errors in
\eqref{eq:shell-error} are precisely radial transport, axial drift mismatch,
and the two cutoff-curvature terms.
\end{remark}

\subsection{Coercive shell control}

\begin{assumption}[Coercive shell control]\label{ass:shell-control}
Let \(u\) be a bounded ancient mild axisymmetric solution with bounded
\(\Gamma\).  Assume in addition that \(u\) is smooth and that the identities
\Cref{lem:gamma-equation,prop:gamma-identity,prop:shell} hold for the weights
below.  There exist expanding smooth weights
\[
   \phi_n(r,z,t)=\lambda_n(r)\eta_n(z-\zeta_n(t)),
   \qquad 0\le\phi_n\le1,\qquad \phi_n\uparrow1
\]
locally uniformly on compact subsets of
\(\mathbb R^3\times(-\infty,0)\), compactly supported in space for each fixed
time.  There are constants \(a_n>0\), \(M_n<\infty\), and
\(\varepsilon_n\ge0\), with \(\varepsilon_n/a_n\to0\), such that
\[
   L_n(t):=\int_{\mathbb R^3}\Gamma^2\phi_n(t)\,dx<\infty
\]
for every \(t<0\), \(L_n\) is locally absolutely continuous, and, for each
\(n\), the following estimates hold for a.e. \(\sigma<0\):
\begin{equation}\label{eq:shell-coercive}
   \int_{\mathbb R^3}|\nabla\Gamma|^2\phi_n(\sigma)\,dx
   \ge
   a_n L_n(\sigma),
\end{equation}
\begin{equation}\label{eq:shell-error-bound}
   \left|
   \int_{\mathbb R^3}\Gamma^2\mathcal E_n[u]\,dx
   \right|
   \le
   2\varepsilon_n,
\end{equation}
and
\begin{equation}\label{eq:shell-past-bound}
   \sup_{\sigma<t}L_n(\sigma)\le M_n .
\end{equation}
Here \(\mathcal E_n[u]\) is the shell-error density associated with
\(\phi_n\).
\end{assumption}

\subsection{Conditional scalar flattening}

\begin{proof}[Proof of \Cref{thm:scalar-flattening-intro}]
Fix \(n\).  Integrating \eqref{eq:shell-identity} from \(s\) to \(t\), with
\(-\infty<s<t<0\), gives
\[
   \frac12L_n(t)-\frac12L_n(s)
   +
   \int_s^t\int_{\mathbb R^3}|\nabla\Gamma|^2\phi_n\,dx\,d\sigma
   =
   \frac12\int_s^t\int_{\mathbb R^3}\Gamma^2\mathcal E_n[u]\,dx\,d\sigma .
\]
Using \eqref{eq:shell-coercive} and \eqref{eq:shell-error-bound}, this implies
\[
   \frac12L_n(t)-\frac12L_n(s)
   +
   a_n\int_s^tL_n(\sigma)\,d\sigma
   \le
   \varepsilon_n(t-s).
\]
Since \(L_n\) is locally absolutely continuous, the preceding integral
inequality yields the differential inequality
\[
   \frac{d}{dt}L_n(t)+2a_nL_n(t)\le2\varepsilon_n
\]
for a.e. \(t<0\).  Multiplying by \(e^{2a_nt}\) and integrating from \(s\) to
\(t\) gives
\[
   L_n(t)
   \le
   e^{-2a_n(t-s)}L_n(s)
   +
   \frac{\varepsilon_n}{a_n}
   \left(1-e^{-2a_n(t-s)}\right).
\]
Using the past bound \eqref{eq:shell-past-bound} and letting
\(s\to-\infty\) yields
\[
   L_n(t)\le \frac{\varepsilon_n}{a_n}.
\]
Let \(K\subset\mathbb R^3\) be compact and fix \(t<0\).  Since
\(\phi_n\uparrow1\) locally uniformly in space-time, for every
\(\eta\in(0,1)\) there is \(n_0=n_0(K,t,\eta)\) such that
\[
   \phi_n(x,t)\ge 1-\eta
   \qquad (x\in K,\ n\ge n_0).
\]
Therefore, for \(n\ge n_0\),
\[
   (1-\eta)\int_K\Gamma(x,t)^2\,dx
   \le
   L_n(t)
   \le
   \frac{\varepsilon_n}{a_n}.
\]
Letting \(n\to\infty\) and using \(\varepsilon_n/a_n\to0\) gives
\[
   \int_K\Gamma(x,t)^2\,dx=0
\]
for every compact \(K\) and every \(t<0\).  Since \(\Gamma\) is continuous by
the smoothness in \Cref{ass:shell-control}, local \(L^2\)-vanishing on every
compact set implies \(\Gamma(x,t)=0\) pointwise for every \(x\) and \(t<0\).
Thus \(\Gamma\equiv0\).
\end{proof}

\subsection{Reduction after scalar flattening}

\begin{corollary}[Reduction after scalar flattening]\label{cor:zero-gamma}
Let \(u\) be as in \Cref{thm:scalar-flattening-intro}.  If in addition \(u\)
is a normalized Type I ancient limit in the same physical frame, then
\(\Gamma\equiv0\), \(u\) belongs to the no-swirl branch, and \(u\equiv0\),
contradicting normalization.
\end{corollary}

\begin{proof}
\Cref{thm:scalar-flattening-intro} gives \(\Gamma\equiv0\).  Then
\(u_\theta=0\) for \(r>0\), so the no-swirl branch closure
\Cref{prop:noswirl-typeI} applies.
\end{proof}

\section{Branch ledger and interface with the residual criterion}
\label{sec:ledger}

\subsection{Closed branches}

\begin{theorem}[Closed structured branch ledger]
\label{thm:closed-structured-ledger}
The following branches cannot contain a nonzero normalized Type I ancient
limit:
\begin{enumerate}[label=\textup{(\roman*)}]
\item axisymmetric no-swirl branch;
\item axisymmetric pointwise \(1/r\) branch;
\item axisymmetric finite-\(L^p\) swirl branch,
\[
   \Gamma\in L_t^\infty L_x^p,\qquad 1\le p<\infty;
\]
\item axisymmetric \(z\)-periodic bounded-\(\Gamma\) branch;
\item perturbative Gallay--Wayne weighted-vorticity branch, whenever the
normalized limit is represented by the associated Biot--Savart velocity in
that branch;
\item bounded-\(\Gamma\) shell-control branch.
\end{enumerate}
\end{theorem}

\begin{proof}
Items (i)--(iv) are \Cref{thm:closed-axisymmetric-branches}.  Item (v) is
\Cref{thm:weighted-vorticity-exclusion-intro}; since the velocity in that
branch is recovered from \(w\) by the Biot--Savart law, \(w\equiv0\) gives the
zero branch representative.  Item (vi) is \Cref{cor:zero-gamma}.  In each
Type I case, the zero conclusion contradicts
\Cref{lem:normalized-nonzero}.
\end{proof}

\subsection{Model branches}

\begin{theorem}[Model branch ledger]
\label{thm:model-branch-ledger}
The Burgers-vortex ancient tube for the strained model consists exactly of the
stationary Burgers vortices \(\alpha G\).  Under the model-reduction assumption
\Cref{ass:bv-reduction}, this rigidity can be transferred to any ancient
solution class mapped into the tube.
\end{theorem}

\begin{proof}
The model rigidity is \Cref{thm:burgers-vortex-tube-intro}; the conditional
transfer is \Cref{cor:bv-import}.
\end{proof}

\subsection{Conditional branches}

\begin{proposition}[Conditional structured exclusions]
\label{prop:conditional-structured}
Under \Cref{hyp:controlled-swirl-liouville}, the controlled-swirl branch
\(\mathfrak C_{\mathrm{cs}}(\gamma)\) contains no nonzero normalized Type I
ancient limit.  Under \Cref{hyp:weighted-velocity-liouville}, the broad
weighted velocity-decay branch \(\mathfrak C_{\mathrm{wd}}(m)\) contains no
nonzero normalized Type I ancient limit.
\end{proposition}

\begin{proof}
This is \Cref{prop:controlled-swirl-conditional,prop:weighted-velocity-conditional}.
\end{proof}

\subsection{Open bounded-swirl obstruction}

\begin{remark}[Open nonperiodic bounded-swirl branch]
The branch not closed unconditionally here is the full class of bounded ancient
axisymmetric unforced Navier--Stokes solutions with nontrivial bounded
\(\Gamma\), no finite-\(L^p\) decay of \(\Gamma\), no axial periodicity, and no
verified shell-control estimate.  The identity \eqref{eq:shell-identity}
reduces this obstruction to explicit transport terms, but it does not
eliminate them.
\end{remark}

\subsection{Structured input for Paper I}

\begin{proof}[Proof of \Cref{thm:structured-input-paperI}]
The closed axisymmetric alternatives are excluded by
\Cref{thm:closed-axisymmetric-branches}.  The perturbative weighted-vorticity
alternative is excluded by \Cref{thm:weighted-vorticity-exclusion-intro} once
the normalized limit is known to lie in that vorticity branch; the
Biot--Savart reconstruction then gives the zero branch representative.  The
bounded-swirl shell branch is excluded by \Cref{cor:zero-gamma}.  These are
the alternatives whose conclusion is zero and hence contradicts the nonzero
normalization of Paper I.  The Burgers-vortex alternative is separate: under
\Cref{ass:bv-reduction}, \Cref{cor:bv-import} transfers the strained-model
conclusion \(\mathcal T[U](\tau)\equiv\bar\alpha G\), and no unforced
zero-solution conclusion is asserted from that model theorem alone.
\end{proof}

\begin{theorem}[Structured/decay input for Paper I]
\label{thm:structured-input}
Let \(\mathcal X_{\mathrm{sd}}\) be the structured/decay branch in Paper I,
restricted to the union of the zero-closure branches in
\Cref{thm:closed-structured-ledger} and the conditional zero-closure branches
for which the corresponding hypotheses above hold.  Then no nonzero normalized
Seregin ancient limit belongs to \(\mathcal X_{\mathrm{sd}}\).  The
Burgers-vortex tube contributes the separate model rigidity statement of
\Cref{thm:model-branch-ledger}.
\end{theorem}

\begin{proof}
This is the zero-closure part of \Cref{thm:structured-input-paperI}, now
phrased as a branch-class statement for Paper I.  The model assertion is
recorded separately in \Cref{thm:model-branch-ledger}.
\end{proof}

\subsection{Dependency ledger}

\begin{center}
\begin{tabular}{p{0.36\textwidth}p{0.18\textwidth}p{0.34\textwidth}}
\hline
Branch & Status & Input \\
\hline
no swirl & closed & KNSS plus Type I constant removal \\
pointwise \(1/r\) & closed & KNSS \\
\(\Gamma\in L_t^\infty L_x^p\) & closed & LZZ plus Type I constant removal \\
periodic bounded \(\Gamma\) & closed & LRZ plus Type I constant removal \\
controlled swirl \(\rho^\gamma V_\theta\) & conditional & \Cref{hyp:controlled-swirl-liouville} \\
broad weighted velocity decay & conditional & \Cref{hyp:weighted-velocity-liouville} \\
perturbative weighted vorticity & closed & Gallay--Wayne Lyapunov functional \\
Burgers-vortex tube & model closed & Gallay--Maekawa Lyapunov functional \\
unforced Burgers-vortex transfer & conditional & \Cref{ass:bv-reduction} \\
bounded nonperiodic \(\Gamma\in L^\infty\) & open; conditional criterion & \Cref{ass:shell-control} \\
bounded \(\Gamma\) plus shell control & closed & scalar flattening plus no-swirl \\
\hline
\end{tabular}
\end{center}

The structured ancient branch is therefore decomposed into several sharply
distinguished mechanisms.  The no-swirl, pointwise scale-invariant,
finite-\(L^p\)-swirl, and periodic bounded-swirl branches are closed by known
axisymmetric Liouville theorems.  The perturbative weighted-vorticity branch is
closed by a Lyapunov functional derived from Gallay--Wayne decay, and
the Burgers-vortex tube is rigid for the strained model by Gallay--Maekawa
stability.  The remaining nonperiodic bounded-swirl branch is not closed here;
it is reduced to the coercive shell-control criterion for the scalar
\(\Gamma=ru_\theta\).  This ledger supplies the structured-branch component of
the residual criterion in Paper I precisely for the closed and conditionally
closed subclasses identified above.

\begin{thebibliography}{99}

\bibitem{CSTY2008}
C.-C. Chen, R. M. Strain, T.-P. Tsai, and H.-T. Yau,
\emph{Lower bound on the blow-up rate of the axisymmetric Navier--Stokes
equations},
Int. Math. Res. Not. IMRN \textbf{2008}, Art. ID rnn016, 31 pp.,
doi:10.1093/imrn/rnn016.

\bibitem{CSTY2009}
C.-C. Chen, R. M. Strain, T.-P. Tsai, and H.-T. Yau,
\emph{Lower bounds on the blow-up rate of the axisymmetric Navier--Stokes
equations II},
Comm. Partial Differential Equations \textbf{34} (2009), no. 3, 203--232,
doi:10.1080/03605300902793956.

\bibitem{DuranPaperI}
G. Duran Ballester,
\emph{Type I blow-up limits and a residual-class criterion for the
three-dimensional Navier--Stokes equations}, preprint, 2026.

\bibitem{DuranPaperII}
G. Duran Ballester,
\emph{Uniformly tight ancient Navier--Stokes dynamics and conditional
equilibrium rigidity}, preprint, 2026.

\bibitem{ESS2003}
L. Escauriaza, G. A. Seregin, and V. \v{S}ver\'ak,
\emph{\(L_{3,\infty}\)-solutions of the Navier--Stokes equations and backward
uniqueness},
Russian Math. Surveys \textbf{58} (2003), no. 2, 211--250,
doi:10.1070/RM2003v058n02ABEH000609.

\bibitem{GallayWayne2002}
T. Gallay and C. E. Wayne,
\emph{Long-time asymptotics of the Navier--Stokes and vorticity equations on
\(\mathbb R^3\)},
Philos. Trans. Roy. Soc. London Ser. A \textbf{360} (2002), no. 1799,
2155--2188, doi:10.1098/rsta.2002.1068.

\bibitem{GallayWayne2006}
T. Gallay and C. E. Wayne,
\emph{Three-dimensional stability of Burgers vortices: the low Reynolds number
case},
Physica D \textbf{213} (2006), no. 2, 164--180,
doi:10.1016/j.physd.2005.11.006.

\bibitem{GallayMaekawa}
T. Gallay and Y. Maekawa,
\emph{Three-dimensional stability of Burgers vortices},
Comm. Math. Phys. \textbf{302} (2011), no. 2, 477--511,
doi:10.1007/s00220-010-1132-6.

\bibitem{Kato1984}
T. Kato,
\emph{Strong \(L^p\)-solutions of the Navier--Stokes equation in
\(\mathbb R^m\), with applications to weak solutions},
Math. Z. \textbf{187} (1984), no. 4, 471--480,
doi:10.1007/BF01174182.

\bibitem{KNSS2009}
G. Koch, N. Nadirashvili, G. Seregin, and V. \v{S}ver\'ak,
\emph{Liouville theorems for the Navier--Stokes equations and applications},
Acta Math. \textbf{203} (2009), no. 1, 83--105,
doi:10.1007/s11511-009-0039-6.

\bibitem{LeiZhang2011}
Z. Lei and Q. S. Zhang,
\emph{A Liouville theorem for the axially-symmetric Navier--Stokes equations},
J. Funct. Anal. \textbf{261} (2011), no. 8, 2323--2345,
doi:10.1016/j.jfa.2011.06.016.

\bibitem{LRZ}
Z. Lei, X. Ren, and Q. S. Zhang,
\emph{A Liouville theorem for axially symmetric Navier--Stokes equations on
\(\mathbb R^2\times\mathbb T^1\)},
Math. Ann. \textbf{383} (2022), no. 1--2, 415--431,
doi:10.1007/s00208-020-02128-9.

\bibitem{LZZ}
Z. Lei, Q. Zhang, and N. Zhao,
\emph{Improved Liouville theorems for axially symmetric Navier--Stokes
equations},
Sci. Sin. Math. \textbf{47} (2017), no. 10, 1183--1198,
doi:10.1360/N012016-00149.

\bibitem{SchmidRossi2004}
P. J. Schmid and M. Rossi,
\emph{Three-dimensional stability of a Burgers vortex},
J. Fluid Mech. \textbf{500} (2004), 103--112,
doi:10.1017/S0022112003007341.

\bibitem{Seregin2012}
G. Seregin,
\emph{A certain necessary condition of potential blow up for Navier--Stokes
equations},
Comm. Math. Phys. \textbf{312} (2012), no. 3, 833--845,
doi:10.1007/s00220-011-1391-x.

\bibitem{SereginSverak2009}
G. Seregin and V. \v{S}ver\'ak,
\emph{On Type I singularities of the local axi-symmetric solutions of the
Navier--Stokes equations},
Comm. Partial Differential Equations \textbf{34} (2009), no. 2, 171--201,
doi:10.1080/03605300802683687.

\bibitem{Sohr2001}
H. Sohr,
\emph{The Navier--Stokes Equations: An Elementary Functional Analytic
Approach},
Birkh\"auser, Basel, 2001,
doi:10.1007/978-3-0348-0551-3.

\end{thebibliography}

\end{document}
